\section{GlyphChain release trust-root governance}
\label{sec:glyphchain-release-trust-root-governance}

Release evidence is meaningful only when the keys allowed to attest it are
themselves authorised through a separate trust process.  Tessaris therefore uses a
three-layer release hierarchy:

\begin{enumerate}
  \item evidence authorities sign task bodies with their individual Ed25519 keys;
  \item a governance authority signs the complete mapping from evidence keys to
        approved roles and independence status; and
  \item the evaluator receives both the governance public key and the exact current
        trust-root inventory digest through an independently controlled channel.
\end{enumerate}

The evaluator does not allow an evidence file to appoint its own verifier, and it
does not treat an unsigned local JSON role assignment as release authority.

\subsection{Versioned signed inventory}

The trust-root document uses schema
\texttt{tessaris.glyphchain-legal-release-evidence-trust-roots.v2}.  Its signed body
contains:

\begin{itemize}
  \item the schema version and UTC generation time;
  \item a map from stable signer identifier to raw Ed25519 public key, approved role
        array and literal boolean independence status; and
  \item the claim boundary explaining that the governance key and current digest
        must be obtained independently.
\end{itemize}

The governance attestation contains its authority identifier, algorithm, canonical
payload digest and Ed25519 signature.  For inventory body $T$ and domain $D_T$,

\begin{equation}
  \sigma_T = \operatorname{Ed25519.Sign}_{sk_G}
  \left(\operatorname{canon}(D_T,T)\right),
\end{equation}

where $D_T$ is
\texttt{TESSARIS\_GLYPHCHAIN\_LEGAL\_RELEASE\_TRUST\_ROOTS\_V1}.  Domain
separation prevents a valid evidence signature or another protocol signature from
being replayed as a trust-root authorization.

\subsection{Independent current-state pin}

A valid signature proves that the governance authority issued an inventory; it does
not by itself prove that the inventory is the latest one.  An older valid inventory
could otherwise be replayed after a signer was removed or had independence revoked.

For exact trust-root file bytes $F_T$, the independently retained current pin is

\begin{equation}
  p_T = \operatorname{SHA256}(F_T).
\end{equation}

Production evaluation of a non-empty inventory requires the operator to supply
$p_T$ separately.  The evaluator hashes the same bounded byte sequence it parses,
compares it with the supplied pin, verifies the governance signature and only then
makes the contained evidence keys available to task evaluation.  Re-reading the
file for later verification is avoided, eliminating a change-between-check-and-use
window.

The public evaluator entry point accepts the bounded raw inventory bytes, the
independently supplied Ed25519 governance public key and the expected file digest.
It derives authentication state and the governance key identifier internally; no
caller-supplied ``authenticated'' boolean is accepted.  The lower-level method
that consumes an already verified root map is private and exists only to separate
the deterministic task engine from input authentication.

The governance public key is also supplied as a separate regular PEM file.  The
readiness report records its key identifier

\begin{equation}
  \mathrm{keyid}_G =
  \texttt{ed25519:sha256:}\operatorname{SHA256}(pk_G)
\end{equation}

and the exact trust-root inventory file digest.  Release operators can therefore
reconstruct which governance anchor and authority set were used for the decision.

\subsection{Whole-inventory validation}

Authentication is necessary but not sufficient.  Before any task can pass, the
entire root map is checked under the following constraints:

\begin{itemize}
  \item no more than 128 evidence authorities;
  \item bounded protocol-safe unique signer identifiers;
  \item exactly one valid 32-byte Ed25519 public key per signer;
  \item no public key appearing under two different signer identifiers;
  \item between one and sixteen unique roles, all drawn from the canonical
        thirty-four-task role vocabulary; and
  \item independence represented only by the literal JSON boolean.
\end{itemize}

The duplicate-public-key rule is essential to separation of duties.  Distinct
labels do not represent distinct authorities when both signatures can be produced
by the same private key.  If any root record is malformed, the report includes a
global inventory error and every supplied evidence task remains failed.

\subsection{Governance inventory builder}

The offline script
\texttt{backend/scripts/create\_glyphchain\_release\_trust\_roots.py} builds the
versioned inventory from evidence-signer public PEM files and explicit approved
role assignments.  It validates every role against the runtime task register,
requires at least one role per signer, applies explicit independence assignments,
sorts roots and roles deterministically and refuses duplicate signer identifiers.

The governance private key must be an owner-only Ed25519 PEM file.  Encrypted PEM
is supported through a named environment variable so the password is not placed in
command history.  The governance key is rejected if it is also one of the evidence
signer keys, preserving the boundary between appointing authorities and supplying
task evidence.

The builder writes two separate exclusive outputs:

\begin{itemize}
  \item the owner-only signed trust-root inventory; and
  \item the public governance key for independent distribution and pinning.
\end{itemize}

It prints the governance key identifier and exact inventory SHA-256 for transfer to
the approved independent channel.  Existing files and symbolic-link destinations
are refused.  If the second output cannot be completed, the newly created inventory
is removed rather than leaving a partial publication set.

\subsection{Evaluation semantics}

Let $S_T$ denote a valid governance signature, $P_T$ an exact current digest match,
$V_T$ complete structural validity and $K_T$ an independently supplied governance
public key.  Trust roots are usable only when

\begin{equation}
  G_T = S_T \land P_T \land V_T \land K_T.
\end{equation}

Task evaluation adds \texttt{trust\_roots\_not\_authenticated} when $S_T$, $P_T$ or
$K_T$ is absent, and \texttt{trust\_root\_inventory\_invalid} when $V_T$ fails.
The report also retains global error details.  Even correctly signed task evidence
cannot pass while $G_T=0$.

The repository's development trust-root inventory is deliberately empty.  An empty
inventory may be evaluated without a governance key only because it authorises no
evidence and necessarily returns a blocked release.  The moment roots are present,
both the independently supplied governance public key and exact current digest are
mandatory.

\subsection{Rotation and revocation procedure}

Removing a signer, changing a role or changing independence produces a new signed
inventory with a new file digest.  Production operators shall publish and pin the
new digest before using it for a release decision.  The old digest must be removed
from the current independent channel while retained in audit history.  Because the
evaluator requires the current exact pin, an old still-valid signature cannot
silently restore a revoked authority.

Governance-key rotation is a separate trust-anchor event.  The replacement public
key and the first inventory it signs must be independently approved and published;
the old private key cannot appoint its own successor merely by embedding it in the
same mutable inventory.  Historical reports retain the prior key identifier and
inventory digest for reproducibility.

\subsection{Regression evidence and claim boundary}

Regressions prove that unsigned non-empty roots cannot close a task, a governance-
signed and independently pinned root set can validate a correctly countersigned
gate, altered role or independence data invalidates the signature, an omitted pin
fails before report creation, the public evaluator cannot be convinced by an
authentication flag, and one public key cannot impersonate two authorities for
separation of duties.

These controls authenticate the configured release authorities.  They do not prove
that an organisation is operationally independent merely because a boolean is set.
That claim still depends on the external governance process, independent custody
and publication evidence required by the controlled-release checklist.
